\scanpage{98}\mspage{42}
% archive sheet number: 97
% L1-4
(2 lisse) Si $X,Y$ sont lisses sur $S$, alors la $W$-équivalence par fibres sur $S$ équivaut à la $W$-équivalence colocale sur $S$, \inserted{et implique la $W$-équiv. \uline{universelle} sur $S$,} donc dans ce cas
\margin{il suffit $=$ $X,Y$ $W$-lisses sur $S$}
% slanted
\diagram{\begin{tikzcd}[ampersand replacement=\&,row sep=tiny,column sep=small]
 b\arrow[r,Leftrightarrow]\arrow[d,Rightarrow]\&c'\\
 c\arrow[d,Rightarrow]\&\\
 a\&
\end{tikzcd}}

% L5-7
(2 \unsure{prop|propr}) Si $X,Y$ sont propres sur $S$, alors la $W$-équiv. par fibres sur $S$ équivaut à la $W$-équivalence locale sur $S$ :
\margin{Il suffit $=$ $X,Y$ $W$-propres sur $S$}
% slanted
\diagram{\begin{tikzcd}[ampersand replacement=\&,row sep=tiny,column sep=small]
 b\arrow[r,Leftrightarrow]\arrow[d,Rightarrow]\&c\\
 c'\arrow[d,Rightarrow]\&\\
 a\&
\end{tikzcd}}

% L8-15
\uline{Dém} 1) n'est autre que L~4 et \deleted{\unsure{L 4 et L 4 bis}} son dual L~4\unsure{$'$|$'{}^{*}$}, pour un localiseur fond. $W$, et on a vu que si on suppose L~1, L~3, alors L~4 équivaut à L~4$'$, et dans la notion de loc. fond. on supposera tjrs que c'est vérifié. Également, nous supposons que L~3 est vraie $=$ sous la forme plus forte L~3 \unsure{quater} (p.~1) : $W_\infty\subset W$, i.e., toute Hot-équivalence \deleted{\unsure{est}}est une $W$-équivalence.
\margin{Mais il y a confusion L~4 est plus faible, cf p.~\unsure{4 §|44|45} l'axiome L~4 bis}
% slanted

% L16-18
(2 \deleted{\unsure{liss}}\unsure{lisse}) Si $X,Y$ lisses sur $S$, alors les $X_s\longrightarrow{}_{s\backslash}X$ $Y_s\longrightarrow{}_{s\backslash}Y$ sont des $W_\infty$-équiv\unsure{,|.} \inserted{donc des $W$-équiv. (il suffirait $=$ que $X,Y$ \deleted{\unsure{soient}} \inserted{\unsure{soient}} $W$-lisses/$S$)} donc le diagramme comm.
\diagram{\begin{tikzcd}[ampersand replacement=\&,row sep=small,column sep=large]
 X_s\arrow[r,"f_s",no head]\arrow[d,"i_s^X"']\&Y_s\arrow[d,"i_s^Y"]\\
 {}_{s\backslash}X\arrow[r,"{}_{s\backslash}f"]\&{}_{s\backslash}Y
\end{tikzcd}}
% L19-24
nous montre que $f_s$ est une $W$-équiv. ssi $ {}_{s\backslash}f$ l'est, d'où $b\Longleftrightarrow c'$. Pour voir que la $W$-équivalence par fibres est \inserted{$=$} une $W$-équivalence \uline{universelle} sur $S$, \deleted{\illegible{1}} on est ramené à montrer que c'est une $W$-équiv. grossière (puisque \deleted{\unsure{universelle}} la lissité stable par chgt. de base sur $S$), \deleted{\unsure{en fait}} \inserted{\unsure{et que|et par cons. que}} \deleted{il devrait suffire $=$ que $X,Y$ soient $W$-lisses sur $S$, \unsure{ou} $W$ \unsure{au} des lisses.} puisque $c'\Rightarrow a$ par 1).

% L25-26
Le cas (2 propre) est dual.

(2 \deleted{\unsure{fibr}}\unsure{fib}) Si $X,Y$ sont des Hot-fibrés sur $S$, alors la
